% !TEX root = ../main.tex
\section{Aggregate Accuracy Does Not Identify Grounding}
\label{sec:independence}

A word on the word. What follows is a \emph{non-identification} result, not a claim of
statistical independence. It shows that an aggregate accuracy figure does not determine,
and without further assumptions does not bound, the grounding rate. It does not show that
the two are probabilistically independent in deployed populations, and they may well be
correlated---a system built to ground its assertions will plausibly be more accurate too.
The distinction matters because the weaker claim is the one a validation regime turns on:
an assessor who accepts an accuracy figure as evidence about grounding is relying on an
entailment that does not hold, whatever the empirical association happens to be.

The redefinition of Section~\ref{sec:hall-def} is not merely a useful
reclassification. It supports a formal result about what aggregate accuracy can and
cannot tell an assessor---and the result is what makes the redefinition consequential
rather than semantic.

\subsection{The result}

For a population of $n$ assertions, define
\begin{equation}
  \mathrm{Accuracy} = \frac{1}{n}\sum_{i=1}^{n} A_i,
  \qquad
  \mathrm{GroundingRate} = \frac{1}{n}\sum_{i=1}^{n} G_i.
  \label{eq:rates}
\end{equation}

\begin{proposition}[Non-identification]
\label{prop:independence}
Let $\alpha, \gamma \in [0,1]$ with $n\alpha$ and $n\gamma$ integral. There is a
population of $n$ assertions with $\mathrm{Accuracy} = \alpha$ and
$\mathrm{GroundingRate} = \gamma$. Aggregate accuracy therefore does not identify the
grounding rate, and supplies no bound on it in the absence of further assumptions; the
same holds in the other direction.
\end{proposition}

\begin{proof}
Let $n_{\mathrm{sound}}, n_{\mathrm{lucky}}, n_{\mathrm{err}}, n_{\mathrm{fab}}$
be the counts in the four cells of Table~\ref{tab:taxonomy}, so that
$n\alpha = n_{\mathrm{sound}} + n_{\mathrm{lucky}}$ and
$n\gamma = n_{\mathrm{sound}} + n_{\mathrm{err}}$. Choose any integer
\[
  n_{\mathrm{sound}} \in
  \bigl[\max(0,\, n\alpha + n\gamma - n),\ \min(n\alpha,\, n\gamma)\bigr],
\]
an interval containing an integer for every admissible $\alpha, \gamma$, and set
$n_{\mathrm{err}} = n\gamma - n_{\mathrm{sound}}$,
$n_{\mathrm{lucky}} = n\alpha - n_{\mathrm{sound}}$, and
$n_{\mathrm{fab}} = n - n\alpha - n\gamma + n_{\mathrm{sound}}$. Each is a
non-negative integer by the choice of interval, and the four sum to $n$.
\end{proof}

The construction requires only that grounding not entail truth, which cell (iii)
supplies: an assertion can rest on authentic, applicable, authoritative evidence
and still be false, through misreading, invalid inference, or evidence adequate for
a weaker claim than the one asserted.

The result is arithmetically simple and strategically significant, and the two
should not be confused. For a population of generated assertions, any feasible
aggregate accuracy can coexist with any feasible grounding rate, because the
population may contain varying proportions of grounded-true, ungrounded-true,
grounded-false, and ungrounded-false assertions. Consequently pooled accuracy
cannot identify or bound the frequency of auditable grounding.

One limitation should be stated before the corollary, because a reader is entitled to
ask for it. The construction establishes \emph{logical} independence: no accuracy
figure entails anything about the grounding rate. It does not establish that the two
are empirically uncorrelated in deployed systems, and they may well be---a system
built to ground its assertions will plausibly be more accurate too. The claim is about
what an accuracy measurement \emph{licenses}, not about what is typical. That is
sufficient for the argument, because a validation regime rests on entailment rather
than on correlation: an assessor who accepts an accuracy figure as evidence of
grounding is relying on a relation that does not hold, whatever the empirical
association happens to be in the population they did not measure.

\begin{corollary}
\label{cor:no-bound}
A system measured at $95\%$ accuracy is consistent with a grounding rate of $0$,
and a system measured at $40\%$ accuracy is consistent with a grounding rate of
$1$.
\end{corollary}

\subsection{What it implies for a validation programme}

What follows is a claim about validation programmes rather than about models. A
programme that reports accuracy without measuring per-assertion grounding cannot
distinguish between radically different reliance profiles---between a system whose
correct answers rest on identifiable authoritative evidence and one whose correct
answers are draws that happened to land well. Those two systems have the same
score and different futures, and the difference appears the moment the question
distribution moves.

Two further readings matter for anyone choosing between releases. Since neither rate constrains the other, an intervention that raises accuracy tells you nothing
about grounding; and optimising against reference answers applies pressure toward
confident guessing, which raises accuracy by moving mass into the lucky cell where
$G_i = 0$. A leaderboard improvement is fully consistent with a system that has
become more willing to assert without support.
